\chapter{Discussion}\label{ch:discussion}
The evidence is most useful when interpreted as a hierarchy of questions with different supervision requirements. This chapter separates direct observations from plausible explanations and the research decisions they motivate.

\section{What the supervision asymmetry explains, and what it does not}
Visible labels accompany repeated photographs, so a surface model can be evaluated throughout the observed visual range. Structural endpoints provide only one measured response per specimen and are embedded in a confounded experimental design. This difference is consistent with stronger visible-condition estimation and weaker structural transfer.

It is tempting to conclude that sparse labels alone explain all structural weakness. The available experiments do not support that decomposition. The saved learning curve varies the number of training specimens, but subset composition changes with it. It does not separate sample size from design coverage or measure all potentially relevant physical factors. Surface appearance may omit information about local metal loss, bond, geometry, or cracking; the selected feature families may also fail to capture available information. The evidence identifies these as unresolved possibilities rather than estimated causes.

The same discipline applies to label adjacency. Near reconstruction explains why a particular total-rust benchmark needs modest interpretation. It does not erase the practical utility of automated surface measurement, and it does not make every peak-rust result tautological. Target-specific evidence should govern the scope of each concern.

\section{Two compatible readings of metadata dominance}
Metadata can carry physically meaningful design information. A mesh configuration and exposure history may legitimately inform a capacity prediction. At the same time, metadata can identify the campaign mixture because those factors co-vary. The pooled results cannot decide how much predictive value belongs to each reading.

Mesh-stratified and campaign-out evaluations narrow what can be claimed. Within a mesh family, a model has less between-group structure to exploit. Under campaign exclusion, it must operate on a combination absent from training. The degraded results limit a claim of transferable specimen-level discrimination. They do not establish that metadata are invalid inputs or that visible corrosion has no physical relationship with capacity.

The baseline diagnostics make this distinction visible. Figure~\ref{fig:ridge-terminal} shows individual terminal errors behind the mean score. Figure~\ref{fig:ridge-coefficients} shows how the full fit allocates coefficient magnitudes among correlated metadata fields without identifying independent effects. Figure~\ref{fig:ridge-learning} shows sensitivity to the number and composition of training specimens within the same mixture. None supplies external validation, calibrated uncertainty or a decomposition of the causes of structural error.

The incremental-value finding should likewise be read narrowly. Image descriptors are tested against a strong metadata baseline for terminal load. The absence of a reliable pooled increment in the saved comparisons is an empirical limit under those inputs, partitions, and selection rules. It is not proof of a universal zero effect; the small post-onset result is one reason to retain that distinction.

\section{Why terminal-load refocusing is scientifically useful}
Terminal load is one of two directly measured structural endpoints. Selecting it as the primary target shortens the validation chain: predictions can be compared with an actual terminal measurement, and image value can be tested against metadata. This permits a sharper question than attempting to validate a succession of inferred hidden states, fitted curves, and crossing times without intermediate ground truth.

The refocus does not turn terminal-image prediction into prognosis. Its primary score is evaluated at the terminal observation, and its training label remains eventual load even for earlier images. A future early-warning study needs a different evaluation design with explicitly restricted observation times. Recognizing this boundary is part of the value of the focused analysis.

\section{Scientific value of the negative findings}
The negative results identify different constraints. The pooled ablation limits the tested image increment. The within-mesh results show how an apparently reasonable MAE can coexist with poor error relative to a narrow target distribution. LOCO exposes extrapolation beyond the observed design combinations. The residual-refinement failure suggests that the tested correction does not resolve the remaining specimen-level gap. The threshold statuses show that a working downstream pipeline need not produce forward crossing estimates.

These findings are mutually informative, but they are not all independent replications. Most share the same specimens; the full image/metadata ablation is centralized in one modelling generation. The report therefore treats agreement as convergence of diagnostics within a dataset, not evidence of broad external validity.

\section{A practical interpretation for research use}
The work can support surface quantification research, comparative evaluation, and decisions about what data should be collected next. It does not support direct capacity certification from a photograph. The distinction follows from the demonstrated evaluation scope, not from a generic objection to machine learning.

The natural next step is to align each desired claim with the measurement and validation needed to test it. Chapter~\ref{ch:limitations} makes the outstanding limitations explicit, and Chapter~\ref{ch:status} converts them into a concrete research agenda.
